\section{Discussion}
\label{sec:discussion}


\paragraph{Headset Optics}
Our study purposefully ignores the impacts of near-eye optics and treats the display stack as an idealized, perfectly head-locked planar display at a fixed VID. This is necessary abstraction in order to generalize our findings to all HMDs, but ignores the realities of real HMD optics. Our simplification treats viewing errors as purely as geometric translations of the view frustum, and this approach isolates the impact of geometric perspective errors while purposefully excluding optical effects. Real viewing errors with near-eye optics are a more complex interaction between the actual displacement from the nominal eye position and the optical design which introduces pupil swim and other optical aberrations. It is important to account for these differences when considering the impacts of viewing errors in real devices.

\paragraph{Implications for IPD Adjustment} Our results indicate that reach behavior remains robust even in the presence of significant IPD errors ($\pm$12 mm). We found that IPD offsets are generally imperceptible except at large magnitudes and close viewing distances (6 mm at 50 mm viewing distance). This aligns with prior research demonstrating a general insensitivity to absolute IPD mismatches in stereoscopic content \cite{allison2015perceptual, didyk2011, guan2016, stevenson1989hyperacuity}. While IPD shifts are known to alter the perceived scale of objects \cite{kim_dwarf_2017, piumsomboon_superman_2018, mine_relationship_2020}, their specific contribution to distance judgment remains unclear. However, perceptual tolerance does not eliminate the need for mechanical IPD adjustment. Physical IPD and IAD alignment is still required to mitigate view-dependent optical artifacts (such as blur and distortion), even if the user is geometrically insensitive to the perspective shift.

\paragraph{Implications for Video Passthrough}
While current commercial HMDs typically rely on static rendering pipelines, our results suggest that the geometric viewing errors inherent to these systems can induce perceptual inaccuracies. This is particularly critical for video passthrough HMDs. Unlike standard VR, where the virtual camera can be placed anywhere, passthrough cameras are physically constrained by the displacement between the external cameras and the user's eyes. Our findings support the adoption of computational imaging methods \cite{kuo23} or view-synthesis algorithms to correct this parallax mismatch \cite{ElChemalyGAP, xiao2022neuralpassthrough}. However, implementation requires a careful trade-off: the latency and potential image artifacts introduced by these methods can sometimes generate registration errors more severe than the original static offset. Consequently, alternative hardware approaches—such as reducing the physical camera-to-eye distance via thinner optics \cite{maimone2020holographic}—may act as a necessary complement to software correction to ensure perceptual accuracy.

\paragraph{Active vs. Static Observers}
Our study focuses on static observation to isolate specific geometric variables; however, real-world HMD usage involves continuous head and eye movement. Perspective projection errors that are imperceptible to a static observer often manifest as dynamic distortions—such as world instability or "swimming"—when the user moves \cite{guan2023perceptual}. These artifacts arise from discrepancies in motion parallax, a depth cue that can supersede stereoscopic cues in sensitivity. Consequently, the detection thresholds reported here may decrease with certain head and eye movements. Future work should extend this protocol to active observers to quantify the true limit at which geometric errors disrupt the stability of the virtual environment.


\section{Conclusion}
We highlighted the perceptual consequences of rendering and viewing errors in HMDs  To do this we first introduced a framework that differentiates between errors in stereo geometry introduced during rendering and viewing stages which can be clearly visualized in an accompanying interactive web application. We next built a software platform that corrects for naturally occurring and unavoidable viewing errors in a Quest 3 HMD which can then be used to purposefully simulate rendering and viewing errors. This platform enabled the set of blind reaching and two-interval forced choice experiments investigating the perceptual consequences of three ecologically-valid rendering and viewing errors in HMDs - a set of experiments that had not been previously possible despite their direct relevance to all egocentric, head-tracked display architectures.

We showed that errors in stereo geometry can induce changes in blind reach behavior but these blind reach distances need to be transformed to correct for reaching bias and put in the egocentric coordinate frame to be interpreted as perceived distance. Distance misperceptions induced by direct passthrough and eye relief errors are well predicted by our geometric model, but depth error induced by IPD error is smaller than predicted by geometry. Evaluating the detectability of IPD error in a 2IFC task revealed that, for a naturalistic and cue-rich scene, even detection thresholds for IPD errors was relatively high at 50 cm, and imperceptible to most viewers at 2.5 m. Overall, our findings demonstrate the perceptual impact of rendering and viewing errors in VR.